\begin{thebibliography}{21}

\bibitem{moore}
Edward F. Moore.
\emph{The Shortest Path Through a Maze}.
In \emph{Proceedings of the International Symposium on the Theory of Switching, Part II}, pp. 285--292. Harvard University Press, 1959.

\bibitem{dijkstra}
Edsger W. Dijkstra.
A note on two problems in connexion with graphs.
\emph{Numerische Mathematik}, 1:269--271, 1959.
\url{https://doi.org/10.1007/BF01386390}

\bibitem{astar}
Peter E. Hart, Nils J. Nilsson, and Bertram Raphael.
A formal basis for the heuristic determination of minimum cost paths.
\emph{IEEE Transactions on Systems Science and Cybernetics}, 4(2):100--107, 1968.
\url{https://doi.org/10.1109/TSSC.1968.300136}

\bibitem{clrs}
Thomas H. Cormen, Charles E. Leiserson, Ronald L. Rivest, and Clifford Stein.
\emph{Introduction to Algorithms}, 4th ed.
MIT Press, 2022.
\url{https://mitpress.mit.edu/9780262046305/introduction-to-algorithms/}

\bibitem{aima}
Stuart Russell and Peter Norvig.
\emph{Artificial Intelligence: A Modern Approach}, 4th ed.
Pearson, 2020.
\url{https://aima.cs.berkeley.edu/}

\bibitem{floyd}
Robert W. Floyd.
Algorithm 97: Shortest path.
\emph{Communications of the ACM}, 5(6):345, 1962.
\url{https://doi.org/10.1145/367766.368168}

\bibitem{warshall}
Stephen Warshall.
A theorem on Boolean matrices.
\emph{Journal of the ACM}, 9(1):11--12, 1962.
\url{https://doi.org/10.1145/321105.321107}

\bibitem{rationalcoding}
Brian Tenneson.
\emph{A Closed-Form Surjection from N onto Q}.
Version 2.2, 2026. Live GitHub copy:
\url{https://github.com/btenneson/pub/blob/main/math.LO_Logic/A%20Closed-Form%20Surjection%20from%20N%20onto%20Q%20by%20Brian%20Tenneson%202.2.pdf}

\bibitem{depths}
Brian Tenneson.
\emph{Depths of the theorem-search program: graph-theoretic proof search, SICs, and automated theorem proving}.
Public 2026 research manuscript in the \emph{Depths} line; live GitHub copy:
\url{https://github.com/btenneson/pub/blob/main/cs.LO_Logic_in_Computer_Science/Depths_of_Induction_ML_Guided_Proof_Search_and_Formalized_Self_Awareness_Version_47.pdf}

\bibitem{horizon}
Brian Tenneson.
\emph{The Universal Proof-Horizon Operator: From MIU and the nth Well-Formed Formula to ZFC, Predator, and Automated Logical Deciders}.
Version 1.0, August 2026. Live GitHub copy:
\url{https://github.com/btenneson/pub/blob/main/cs.LO_Logic_in_Computer_Science/The_Universal_Proof-Horizon_Operator_by_Brian_Tenneson.pdf}

\bibitem{data201}
Brian Tenneson.
\emph{DATA 2.0.1: Proof-Horizon Architecture / From a Target Formula to a Verified Minimum-Cost Proof}.
August 2026. Live GitHub copy:
\url{https://github.com/btenneson/pub/blob/main/cs.LO_Logic_in_Computer_Science/DATA_2_0_1_Proof_Horizon_Architecture_FINAL.pdf}

\bibitem{horizoncomp}
Brian Tenneson.
\emph{A Technical Companion to My Next Generation Automated Theorem Prover}.
2026; includes the Proof-Horizon/Ocean synthetic benchmark material used here. Live GitHub copy:
\url{https://github.com/btenneson/pub/blob/main/cs.LO_Logic_in_Computer_Science/A_Technical_Companion_to_My_Next_Generation_Automated_Theorem_Prover_v1.0.pdf}

\bibitem{hilbert}
Brian Tenneson.
\emph{Hilbert Space-Filling Curve Theorem Search}.
Version 0.5.3, 2026. Live GitHub copy:
\url{https://github.com/btenneson/pub/blob/main/cs.LO_Logic_in_Computer_Science/Hilbert_Space_Filling_Curve_Theorem_Search_v0.5.3.pdf}

\bibitem{data3}
Brian Tenneson.
\emph{DATA 3.0: AMLD/IFS Semigroup Architecture, Settlement Horizons, and Conjecture-Settling Research Program}.
August 2026. Live GitHub research library:
\url{https://github.com/btenneson/pub/tree/main/cs.LO_Logic_in_Computer_Science}

\bibitem{godel}
Kurt Goedel.
Ueber formal unentscheidbare Saetze der Principia Mathematica und verwandter Systeme I.
\emph{Monatshefte fuer Mathematik und Physik}, 38:173--198, 1931.
\url{https://doi.org/10.1007/BF01700692}

\bibitem{church}
Alonzo Church.
An unsolvable problem of elementary number theory.
\emph{American Journal of Mathematics}, 58(2):345--363, 1936.
\url{https://doi.org/10.2307/2371045}

\bibitem{tarski}
Alfred Tarski.
Der Wahrheitsbegriff in den formalisierten Sprachen.
\emph{Studia Philosophica}, 1:261--405, 1935.

\end{thebibliography}

\end{document}
